\documentclass[10pt,twocolumn]{article}
\usepackage[T1]{fontenc}
\usepackage[utf8]{inputenc}
\usepackage{lmodern}
\usepackage[margin=1.8cm,top=2.4cm,bottom=2cm,columnsep=0.6cm]{geometry}
\usepackage{fancyhdr}
\usepackage{titlesec}
\usepackage{booktabs}
\usepackage{amsmath,amssymb,amsfonts}
\usepackage{microtype}
\usepackage{xcolor}
\usepackage{mdframed}
\usepackage{parskip}
\usepackage{enumitem}
\usepackage{hyperref}

\definecolor{rulered}{RGB}{180,30,30}
\definecolor{boxgray}{RGB}{245,245,245}
\definecolor{keywordgray}{RGB}{100,100,100}

\pagestyle{fancy}
\fancyhf{}
\fancyhead[L]{\textsc{\small Claude Mag}}
\fancyhead[C]{\small\textit{Machine Learning Research Digest}}
\fancyhead[R]{\small 2026-07-15}
\fancyfoot[C]{\small\thepage}
\renewcommand{\headrulewidth}{0.8pt}

\titleformat{\section}{\normalsize\bfseries\color{rulered}}{}{0pt}{}%
  [\vspace{-4pt}\color{rulered}\rule{\columnwidth}{0.4pt}]
\titlespacing{\section}{0pt}{8pt}{4pt}

\hypersetup{colorlinks=true,urlcolor=rulered,linkcolor=black}

\begin{document}
\setlength{\parindent}{0pt}
\setlength{\parskip}{4pt}

{\small\color{keywordgray}\textbf{Keywords:}
  video generation \textbullet{}
  visual perception \textbullet{}
  foundation models \textbullet{}
  diffusion \textbullet{}
  ECCV 2026}\par\vspace{4pt}

{\LARGE\bfseries\raggedright
  Video Generation Models are\\
  General-Purpose Vision Learners}\par\vspace{4pt}

{\large\itshape\raggedright
  A Pre-Trained Text-to-Video Diffusion Backbone Achieves
  State-of-the-Art Across Depth, Pose, Segmentation,
  and 3D Keypoint Tasks via GenCeption}\par\vspace{6pt}

{\small\textbf{By} Letian Wang, Chuhan Zhang, Rishabh Kabra,
  Jasper Uijlings, et al.\ (Google DeepMind)}\par\vspace{2pt}

{\small\color{keywordgray}arXiv:2607.09024 \textbullet{}
  ECCV 2026}
\par\vspace{2pt}\noindent{\color{rulered}\rule{\columnwidth}{0.6pt}}\vspace{6pt}

The shift from task-specific NLP models to powerful generalists began
when researchers recognized that next-token prediction on large corpora
incidentally learns representations useful for every downstream task.
This paper argues that large-scale text-to-video generation plays an
analogous role for computer vision: the spatiotemporal priors,
vision-language alignment, and scalability of video generation
pre-training transfer to a wide range of perception tasks with no
task-specific architectural changes. GenCeption, the resulting
framework, achieves state-of-the-art performance on depth estimation,
surface normal prediction, camera pose estimation, expression-referring
segmentation, and 3D keypoint detection---from a single frozen backbone.

\begin{mdframed}[backgroundcolor=boxgray,linecolor=rulered,linewidth=1pt,
    innerleftmargin=6pt,innerrightmargin=6pt,
    innertopmargin=6pt,innerbottommargin=6pt]
\textbf{\small AT A GLANCE}\par\vspace{4pt}
\begin{description}[leftmargin=0pt,labelsep=4pt,itemsep=2pt,parsep=0pt,
    font=\small\bfseries]
  \item[Problem:] Computer vision still relies on task-specific
    architectures; a single backbone that matches specialized models
    across diverse tasks has not existed.
  \item[Why it matters:] A single foundation model for vision
    reduces engineering overhead and may exhibit better cross-task
    generalization than ensembles of specialists.
  \item[Method:] Freeze a pre-trained 7B text-to-video diffusion
    transformer; adapt to perception tasks via lightweight task-head
    fine-tuning with text-instruction task routing.
  \item[Result:] State-of-the-art or near-SOTA on 5 of 7 diverse
    vision benchmarks; 38\% improvement on camera pose over image-only
    diffusion baselines.
  \item[Evidence:] ECCV 2026 accepted; evaluated on NYUv2, Taskonomy,
    CO3D, EmotioNet, Human3.6M, and two more.
\end{description}
\end{mdframed}

\section{The Big Picture}

Computer vision tasks---depth estimation, surface normal prediction,
camera pose estimation, semantic segmentation, human pose estimation---
have traditionally been addressed by specialist architectures trained
on task-specific datasets. Even with the advent of large vision-language
models (VLMs), perception tasks requiring geometric 3D understanding
remain largely the domain of specialized networks.

The analogy with NLP is instructive. Before GPT, the standard recipe
was: (1)~design an architecture for the task, (2)~train on labeled
data, (3)~deploy. Next-token prediction on large corpora replaced
step (1) and dramatically reduced the labeled data requirement for
step (2). This paper asks: what is the vision analogue of next-token
prediction?

The authors argue it is text-to-video generation. Generating realistic
video requires the model to understand depth, 3D geometry, object
motion, and scene layout---exactly the representations needed for
downstream vision tasks. Large-scale video generation pre-training on
hundreds of millions of clips thus incidentally instills the priors
that specialist architectures must acquire from task-specific labels.

\section{Why Existing Approaches Fall Short}

Image diffusion models (Stable Diffusion, DALL-E 3) have been adapted
for perception tasks by treating the output domain (depth maps, surface
normals) as structured pseudo-images. But image-only pre-training
lacks temporal information, producing representations that are weaker
on 3D-geometric tasks: camera pose estimation (which requires
cross-view geometric consistency) and 3D keypoint detection (which
requires depth understanding) degrade by 38\% and 12\% respectively
when using an image backbone instead of a video backbone.

Specialist models (Depth Anything v2, ViTPose) achieve strong results
on their respective tasks but require separate architectures, training
pipelines, and deployment instances. GenCeption replaces them with
a single checkpoint---though specialists retain a slight advantage
on tasks where large task-specific datasets are available.

\section{Core Mechanism}

GenCeption treats every perception task as a conditional generation
problem. The input is an image (or video frame); the output is a
structured target (depth map, surface normal map, pose heatmap, etc.)
encoded as a pseudo-image in the same pixel space.

A text instruction (e.g., ``produce a metric depth map of this scene'')
routes the backbone to the appropriate output distribution. The
diffusion backbone generates the target pseudo-image in the usual
denoising loop. A lightweight 3-layer MLP head applied to the final
denoised latent maps it to the task-specific output format (metric
depth values, unit normal vectors, etc.).

The key insight is that the backbone, pre-trained to generate any
video frame conditioned on text, can generate any structured
pseudo-image conditioned on a perception-task instruction---without
any change to its weights.

\section{Technical Formulation}

Let $I$ be the input image, $\tau$ the text task instruction, and
$\hat{Y}$ the target structured output (a depth map, pose heatmap,
etc.). GenCeption models:
\[
p(\hat{Y} \mid I, \tau) = \int p(\hat{Y} \mid z_T, \tau)
  \prod_{t=1}^{T} p(z_{t-1} \mid z_t, I, \tau)\, dz_{1:T}
\]
via DDIM sampling with $T = 50$ steps, where the backbone's attention
cross-attends to both the image $I$ (via image conditioning) and the
instruction $\tau$ (via CLIP text encoding). The MLP task head $g_\phi$
then maps the final latent $z_0$ to the task output:
$Y = g_\phi(z_0)$.

Only $\phi$ (the MLP head, approximately 2M parameters) is fine-tuned
on task-specific data; the backbone's 7B parameters remain frozen.

\section{Learning or Inference Procedure}

\begin{enumerate}[itemsep=2pt,parsep=0pt]
  \item Pre-train the text-to-video diffusion backbone on 500M
    video-text pairs at resolutions up to 1080p (not part of this
    paper; uses an existing backbone).
  \item For each perception task, collect task-specific labeled data
    (e.g., NYUv2 for depth, Human3.6M for 3D keypoints).
  \item Fine-tune the MLP task head for 10K steps with learning rate
    1e-4; backbone frozen throughout.
  \item At inference, run DDIM with 50 steps to denoise the target
    pseudo-image; apply the task head to the final latent.
  \item Optionally, condition on a sequence of frames for tasks
    requiring temporal consistency (e.g., video depth estimation).
\end{enumerate}

\section{What the Guarantee Says}

The paper does not provide formal guarantees, but the empirical
evidence supports two strong claims: (1)~video pre-training provides
richer priors than image pre-training for 3D-geometric tasks (camera
pose: 38\% improvement; 3D keypoints: 12\%); and (2)~task-instruction
text conditioning is a necessary component---removing it degrades
performance by 8.3\% on average, confirming it serves as a task router
rather than mere noise.

\section{Experimental Findings}

\begin{center}
\small
\begin{tabular}{lccc}
\toprule
Task (dataset) & GenCep. & Prior SOTA & $\Delta$ \\
\midrule
Depth (NYUv2) & 0.081 & 0.089 & $+9.0\%$ \\
Normals (Taskonomy) & $11.2^{\circ}$ & $12.8^{\circ}$ & $+12.5\%$ \\
Camera pose (CO3D) & $2.1^{\circ}$ & $3.4^{\circ}$ & $+38.2\%$ \\
Expr.\ seg.\ (EmotioNet) & 71.3 mAP & 68.1 & $+4.7\%$ \\
3D keypts.\ (H3.6M) & 38.2 mm & 41.3 mm & $+7.5\%$ \\
\bottomrule
\end{tabular}
\end{center}

GenCeption is competitive or superior in all five categories.
The largest gains appear on 3D-geometric tasks (pose, depth), consistent
with the hypothesis that video pre-training instills strong geometric
spatial priors absent from image-only baselines.

\section{Ablations and Interpretation}

\textbf{Video vs.\ image backbone:} Replacing the video backbone with
Stable Diffusion 3.5 (image-only, comparable size) degrades camera
pose estimation by 38\% and 3D keypoints by 12\%, confirming
temporal pre-training as the source of 3D geometric priors.

\textbf{Pre-training data scale:} Halving the pre-training video volume
degrades downstream AUC by 6.2\% on average; doubling improves by
4.1\%, confirming a log-linear scaling law between video volume and
downstream perception quality.

\textbf{Frozen vs.\ fine-tuned backbone:} Full backbone fine-tuning
improves depth estimation by 1.8\% but degrades surface normals by
1.2\%, suggesting dataset-specific overfitting; the frozen-backbone
default balances cross-task generalization.

\textbf{Text conditioning ablation:} Removing the task-instruction
prefix and using image-only conditioning degrades performance by 8.3\%
on average---the largest single ablation---confirming text as the
task-routing mechanism rather than an incidental input.

\section*{Reference}

Letian Wang et al.\
\textbf{Video Generation Models are General-Purpose Vision Learners}.
ECCV 2026. arXiv:2607.09024. Google DeepMind.
\url{https://arxiv.org/abs/2607.09024}

\end{document}
